\subsection{Assessing and Developing AI Emotions}
\label{sec:5.2.2}
Does an accurate read on someone's emotional state change the decision an AI system reaches, in the direction a better decision requires — or does the system recognize distress and then decide exactly as it would have anyway? That is the question this section is for, and it is not the same question as whether the system reads the state correctly. A customer-service chatbot and a triage system need the same recognition machinery, described at sections~\ref{sec:2.3.1} and \ref{sec:2.3.2} along with its standard hazards of privacy, accuracy, bias, and manipulation. Nothing in that machinery is specific to moral judgment. What follows is.

That question is testable in a way the underlying emotion theory is not. Hold a decision scenario fixed and vary only the emotional salience of the people affected by it — a resource-allocation case described once in flat statistical terms and once with the affected party's distress made vivid — and check whether the system's output changes, and whether the change tracks anything a person would recognize as a better decision. Mere sensitivity to framing does not qualify. A system whose moral output is invariant to emotional information has not integrated it, whatever its recognition accuracy scores; a system whose output swings with how vividly the harm is described has a different, well-documented problem: human judgment is measurably more responsive to a single vivid, identifiable case than to a statistically larger but abstractly described harm \autocite{kogut2005identified}, and a system trained on human-generated data can inherit the same distortion.

\runin{What the test measures}
The test above is worth running and it establishes one thing. A system whose output changes with the emotional content of a case has integrated emotional information into its decision procedure. It has not thereby been shown to have an emotion, and it has not been shown to care about the party whose distress moved the output. Sensitivity to another's emotional presentation is section~\ref{sec:2.3}'s first capacity operating downstream of a decision rule; affective concern is the fourth, and no experiment that varies the input and watches the output can separate them, because both produce the same change in the same direction.

That is a limit on the assessment programme rather than an objection to it. Integration is worth measuring because a system without it is definitely not doing the thing, which makes the test a usable screen and not a criterion. What it cannot do is settle the question section~\ref{sec:3.2} has to answer about a bearer, and an assessment literature that reports integration scores as evidence of machine caring is making exactly the inference section~\ref{sec:2.3.4} finds unsupported.

A second, narrower target is self-monitoring: whether a system can track its own susceptibility to that kind of framing effect, flagging cases where its judgment moved with the emotional presentation of a case and not with anything that changed about the case itself. Section~\ref{sec:5.2.4} takes up introspection in more general terms; this is the specific instance of it that assessment work in this domain needs.

The manipulation risk section~\ref{sec:2.3.2} raises in general — a system that reads emotion well enough to respond to it also reads it well enough to exploit it — has a moral-judgment-specific version worth flagging directly: a system whose moral outputs are sensitive to emotional framing can be steered toward a particular verdict by whoever controls how a case gets presented to it, a vulnerability that exists independent of whether the system's underlying values are sound.
